\newpage

\section{মুজারে}

% ─── Fil Mudare Table ───────────────────────────────────────
\begin{center}
  {\arabicfont \color{titlegreen} \textbf{\textarabic{المضارع المعروف}}}
\end{center}

\vspace{0.3cm}

\begin{center}
  {\arabicfont
    \textarabic{%
      \textbf{\color{titlegreen}علامة:}
      \enspace
      الجذر \textbf{\myroot{فعل}}
      \enspace|\enspace
      حرف المضارعة والعلامة \textbf{\color{suffixred}بالأحمر}%
    }
  }
\end{center}

\vspace{0.4cm}

\setlength{\arrayrulewidth}{0.5pt}
\setlength{\tabcolsep}{3pt}
\renewcommand{\arraystretch}{2.2}

\noindent\begin{tabularx}{\textwidth}{|C|C|C|}

  \hline
  \rowcolor{headergreen}
  {\arabicfont \color{white} \textbf{مُتَكَلِّم}} &
  {\arabicfont \color{white} \textbf{حَاضِر}}     &
  {\arabicfont \color{white} \textbf{غَائِب}}     \\
  \hline

  % Row 1 — واحد مذكر
  \rowcolor{lightgreen}
  {\arabicfont \mudraf{أَ}}            &
  {\arabicfont \mudraf{تَ}}            &
  {\arabicfont \mudraf{يَ}}            \\
  \hline

  % Row 2 — تثنية مذكر / جمع متكلم
  {\arabicfont \mudraf{نَ}}            &
  {\arabicfont \mud{تَ}{فْعَلَ}{انِ}}  &
  {\arabicfont \mud{يَ}{فْعَلَ}{انِ}}  \\
  \hline

  % Row 3 — جمع مذكر
  \rowcolor{lightgreen}
  &
  {\arabicfont \mud{تَ}{فْعَلُ}{ونَ}}  &
  {\arabicfont \mud{يَ}{فْعَلُ}{ونَ}}  \\
  \hline

  % Row 4 — واحد مؤنث
  &
  {\arabicfont \mud{تَ}{فْعَلِ}{ينَ}}  &
  {\arabicfont \mudraf{تَ}}            \\
  \hline

  % Row 5 — تثنية مؤنث
  \rowcolor{lightgreen}
  &
  {\arabicfont \mud{تَ}{فْعَلَ}{انِ}}  &
  {\arabicfont \mud{تَ}{فْعَلَ}{انِ}}  \\
  \hline

  % Row 6 — جمع مؤنث
  &
  {\arabicfont \mud{تَ}{فْعَلْ}{نَ}}   &
  {\arabicfont \mud{يَ}{فْعَلْ}{نَ}}   \\
  \hline

\end{tabularx}

\vspace{0.4cm}

\begin{center}
  {\arabicfont \color{gray} \small
    \textarabic{حرف المضارعة (يَ تَ أَ نَ) والعلامة بالأحمر \textbullet\ الجذر بالأسود}
  }
\end{center}

\vspace{0.45cm}

\noindent মাজহুল মুজারে {\arabicfont \majhulmud}:
হরফে মুজারায় পেশ, আইনে যবর।

\vspace{0.35cm}

\noindent মুজারেতে চারটি কণা বসে।
পাঁচ সিগায় শেষ হরকত বদলায়: গায়েবের পুং ও স্ত্রী একবচন, হাজিরের পুং একবচন, মুতাকাল্লিমের দুই সিগা।
সাত সিগা থেকে নুনে ই'রাব পড়ে। স্ত্রী বহুবচনের দুই সিগা আগের মতো থাকে।

\vspace{0.25cm}

\begingroup
\setlength{\arrayrulewidth}{0.5pt}
\setlength{\tabcolsep}{3pt}
\renewcommand{\arraystretch}{1.8}

\noindent\begin{tabularx}{\textwidth}{|C|C|C|}
  \hline
  \rowcolor{headergreen}
  {\color{white}\textbf{কণা}} &
  {\color{white}\textbf{রূপ}} &
  {\color{white}\textbf{ফল}} \\
  \hline
  \rowcolor{lightgreen}
  {\arabicfont \textarabic{لَا}} &
  {\arabicfont \mudla} &
  না-বাচক, শেষে পেশ \\
  \hline
  {\arabicfont \textarabic{لَنْ}} &
  {\arabicfont \mudlan} &
  দৃঢ় না, শেষে যবর \\
  \hline
  \rowcolor{lightgreen}
  {\arabicfont \textarabic{لَمْ}} &
  {\arabicfont \mudlam} &
  অতীতের না, শেষে সাকিন \\
  \hline
  {\arabicfont \textarabic{لَـ نَّ}} &
  {\arabicfont \mudnun} &
  নিশ্চয়তা, শেষে যবর \\
  \hline
\end{tabularx}
\endgroup

\vspace{0.25cm}

\noindent ভারী নুন {\arabicfont \textarabic{نَّ}} একবচন ও পুং বহুবচনে।
হালকা নুন {\arabicfont \textarabic{نْ}} দ্বিবচন ও স্ত্রী বহুবচনে; তখন আগে আলিফ বসে।
লাম ও লান মুজারেকে ভবিষ্যৎ কালে বেঁধে দেয়।
